% 第16章 SU(N) Heisenberg 模型
\chapter{$SU(N)$ Heisenberg 模型}

过去十年间，用大 $N$ 近似处理强相互作用量子体系的工作极为广泛。这一途径起源于基本粒子理论，但在凝聚态物理中找到了众多应用（见文献注释）。最初，大 $N$ 展开是为金属中磁性杂质的 Kondo 与 Anderson 模型发展的；不久之后扩展到稀土化合物中混合价与重费米子现象的 Kondo 与 Anderson 格子模型；也被用于处理高 $T_{c}$ 铜氧化物超导体中的强 Coulomb 相互作用。

在本章及第 17、18 章中，我们将把大 $N$ 途径形式化并应用于量子 Heisenberg 模型。这一方法为量子磁体的静态与动力学关联提供了又一条途径。下面导出的平均场理论既能描述有序相也能描述无序相，零温与有限温皆可，它们与第 11--15 章的半经典途径互为补充。

一般地说，参数 $N$ 标记每个格点上内部 $SU(N)$ 对称的“味”数（即一个 Schwinger 玻色子或受约束费米子可取的味的数目）。多数情况下，大 $N$ 近似被用于处理对称性为 $SU(2)$ 的自旋哈密顿量，因此 $N$ 并不是真正大的参数。尽管如此，$1/N$ 展开提供了得到简单平均场理论的便捷方法。实践发现这些理论或者出奇地成功，或者完全错误，取决于具体体系。例如第 18 章将看到：一维的 Schwinger 玻色子平均场理论对铁磁体和整数自旋反铁磁体行之有效，但对半奇整数自旋反铁磁体则失效。

大 $N$ 途径用约束处理强的局域相互作用。它不是按相互作用大小的微扰展开，而是鞍点展开，通常保持哈密顿量的自旋对称。平均场水平上，约束只在平均意义上被强制；其效应由 $1/N$ 的更高阶修正系统地重新引入。

平均场理论的 $1/N$ 修正可以用 Feynman 图计算，后者描述平均场理论自由准粒子之间的相互作用。17.3 节我们将利用 $1/N$ 展开的独特结构，导出对一切阶成立的某些求和规则。

事实表明，不同的大 $N$ 推广适用于不同的 Heisenberg 模型，取决于耦合符号、自旋大小与格子。“适用”指相应的平均场理论给出基态正确的定性特征。下面引入 Heisenberg 模型的三种大 $N$ 推广，从铁磁体的 Schwinger 玻色子表示开始。

\section{铁磁体：Schwinger 玻色子}

循 7.2 节，每格点引入两个 Schwinger 玻色子 $a_{i}, b_{i}$，并在其 Fock 空间上施加局域约束

\begin{equation*}
a _ {i} ^ {\dagger} a _ {i} + b _ {i} ^ {\dagger} b _ {i} = 2 S.
\tag{16.1}
\end{equation*}

定义键算符

\begin{equation*}
\mathcal {F} _ {ij} \equiv a _ {i} ^ {\dagger} a _ {j} + b _ {i} ^ {\dagger} b _ {j}.
\tag{16.2}
\end{equation*}

自旋 $S$ 的铁磁 Heisenberg 模型写为

\begin{align*}
\mathcal {H} & = - J \sum_{\langle ij \rangle} \mathbf {S} _ {i} \cdot \mathbf {S} _ {j}\\
& = - \frac {J}{2} \sum_{\langle ij \rangle} \left[ \mathcal {F} _ {ij} ^ {\dagger} \mathcal {F} _ {ij} - 2 S ( S + 1 ) \right]\\
& = - \frac {J}{2} \sum_{\langle ij \rangle} \left( : \mathcal {F} _ {ij} ^ {\dagger} \mathcal {F} _ {ij} : - 2 S ^ {2} \right) .
\tag{16.3}
\end{align*}

$\sum_{\langle ij\rangle}$ 对以 $(i,j)$ 为端点的键求和，$: :$ 表示正规排序（见 (C.18)）。把 Schwinger 玻色子的味数从 2 增到 $N$，(16.3) 推广到 $SU(N)$：

\begin{equation*}
\left( a _ {i}, b _ {i} \right) \rightarrow \left( a _ {i 1}, a _ {i 2}, \dots , a _ {iN} \right).
\tag{16.4}
\end{equation*}

键算符推广为

\begin{equation*}
\mathcal {F} _ {ij} \rightarrow \sum_ {m = 1} ^ {N} a _ {im} ^ {\dagger} a _ {jm} ,
\tag{16.5}
\end{equation*}

约束为

\begin{equation*}
\sum_ {m = 1} ^ {N} a _ {im} ^ {\dagger} a _ {im} = NS \, , \quad S = 0, 1, 2, \ldots .
\tag{16.6}
\end{equation*}

$SU(N)$ 铁磁玻色子（FM-B）Heisenberg 模型为

\begin{align*}
\mathcal {H} ^ {\mathrm{FM-B}} ( N ) & = - \frac {J}{N} \sum_{\langle ij \rangle} \left( : \mathcal {F} _ {ij} ^ {\dagger} \mathcal {F} _ {ij} : - NS ^ {2} \right)\\
& = - \frac {J}{N} \sum_ {ij} \left( \sum_ {mm '} S _ {i} ^ {mm '} S _ {j} ^ {m'm} - NS ^ {2} \right) ,
\tag{16.7}
\end{align*}

第二行由重排算符并把 $SU(N)$ 生成元定义为

\begin{equation*}
S ^ {mm '} = a _ {m} ^ {\dagger} a _ {m '}.
\tag{16.8}
\end{equation*}

得到。$S^{mm'}$ 满足 $\mathrm{SU}(N)$ 代数

\begin{equation*}
\left[ S ^ {mm '}, S ^ {\mu \mu '} \right] = \delta_{m' \mu} S^{m \mu '} - \delta_{m \mu '} S^{\mu m'}.
\tag{16.9}
\end{equation*}

特定的“自旋大小” $S$ 由约束 (16.6) 定义。对 $SU(2)$，我们认出通常的自旋算符：$S^{12} = S^{+}$，$S^{21} = S^{-}$，$S^{11} = S^{z} + S$。

可以验证 $H^{\mathrm{FM-B}}$ 在自旋的均匀 $SU(N)$ 变换下不变。我们有意把哈密顿量写成负二次型 $-F_{ij}^{\dagger}F_{ij}$；如将见，这为平均场理论中解耦相互作用提供了自然的出发点。

\section{反铁磁体：Schwinger 玻色子}

考虑子晶格 A、B 的二分格子上近邻反铁磁相互作用 $J > 0$ 的情形。键 $\langle ij\rangle$ 定义为 $i \in A$、$j \in B$。反铁磁键算符定义为

\begin{equation*}
\vec {\mathcal {A}} _ {ij} = a _ {i} b _ {j} - b _ {i} a _ {j}.
\tag{16.10}
\end{equation*}

箭头 $\rightarrow$ 表示对 $i \rightarrow j$ 交换的反对称。在子晶格 B 上定义绕 $y$ 轴转 $\pi$ 的自旋转动：

\begin{equation*}
a _ {j} \rightarrow - b _ {j} \, , \quad b _ {j} \rightarrow a _ {j}.
\tag{16.11}
\end{equation*}

这是保持约束 (16.1) 的正则变换。反铁磁键算符变成对称算符：

\begin{equation*}
\vec {\mathcal {A}} _ {ij} \rightarrow \mathcal {A} _ {ij} = a _ {i} a _ {j} + b _ {i} b _ {j}.
\tag{16.12}
\end{equation*}

$SU(2)$ Heisenberg 模型写为

\begin{align*}
\mathcal {H} & = J \sum_{\langle ij \rangle} \mathbf {S} _ {i} \cdot \mathbf {S} _ {j}\\
& = - \frac {J}{2} \sum_{\langle ij \rangle} \left( \mathcal {A} _ {ij} ^ {\dagger} \mathcal {A} _ {ij} - 2 S ^ {2} \right).
\tag{16.13}
\end{align*}

与铁磁情形一样，可以通过增加 Schwinger 玻色子的味把 $\mathcal{H}$ 推广到 $N > 2$ 模型。约束推广为 (16.6)，键算符推广为

\begin{equation*}
\mathcal {A} _ {ij} \rightarrow \sum_ {m = 1} ^ {N} a _ {im} a _ {jm}.
\tag{16.14}
\end{equation*}

$SU(N)$ 反铁磁玻色子（AFM-B）Heisenberg 模型为

\begin{align*}
\mathcal {H} ^ {\mathrm{AFM-B}} ( N ) & = - \frac {J}{N} \sum_{\langle ij \rangle} \left( \mathcal {A} _ {ij} ^ {\dagger} \mathcal {A} _ {ij} - NS ^ {2} \right)\\
& = - \frac {J}{N} \sum_{\langle ij \rangle} \left( \sum_ {mm '} S _ {i} ^ {mm '} \tilde {S} _ {j} ^ {m'm} - NS ^ {2} \right) ,
\tag{16.15}
\end{align*}

其中

\begin{equation*}
\tilde {S} _ {j} ^ {mm '} = a _ {jm '} ^ {\dagger} a _ {jm}
\tag{16.16}
\end{equation*}

是子晶格 B 上共轭表示的生成元。注意 (16.15) 的 $H^{\mathrm{AFM-B}}$ 在均匀 $SU(N)$ 变换 $U$ 下不变，只在子晶格 A、B 上分别作交错共轭转动 $U$ 与 $U^{\dagger}$ 下不变。

\section{反铁磁体：受约束费米子}

第 2 章我们首次遇到作为电子双线性算符的量子自旋。自旋 1/2 算符可以用每格点两个态 $|i,\uparrow\rangle$、$|i,\downarrow\rangle$ 表示，各被一个费米子占据：

\begin{equation*}
n _ {i \uparrow} + n _ {i \downarrow} = 1 \; ( = 2 S ).
\tag{16.17}
\end{equation*}

由于 Pauli 原理，这一表示限于 $S = \frac{1}{2}$\footnote{$S > \frac{1}{2}$ 的表示可以通过增加轨道与附加约束构造。}。自旋算符为（见 (2.10)）

\begin{align*}
S _ {i} ^ {+} & = c _ {i \uparrow} ^ {\dagger} c _ {i \downarrow} ,\\
S _ {i} ^ {-} & = c _ {i \downarrow} ^ {\dagger} c _ {i \uparrow} ,\\
S _ {i} ^ {z} & = \frac {1}{2} \left( c _ {i \uparrow} ^ {\dagger} c _ {i \uparrow} - c _ {i \downarrow} ^ {\dagger} c _ {i \downarrow} \right) .
\tag{16.18}
\end{align*}

键算符定义为

\begin{equation*}
\mathcal {D} _ {ij} = c _ {i \uparrow} ^ {\dagger} c _ {j \uparrow} + c _ {i \downarrow} ^ {\dagger} c _ {j \downarrow} ,
\tag{16.19}
\end{equation*}

Heisenberg 模型写为

\begin{align*}
\mathcal {H} & = J \sum_{\langle ij \rangle} \mathbf {S} _ {i} \cdot \mathbf {S} _ {j}\\
& = - \frac {J}{2} \sum_{\langle ij \rangle} : \mathcal {D} _ {ij} ^ {\dagger} \mathcal {D} _ {ij} : .
\tag{16.20}
\end{align*}

引入 $N$ 味费米子 $c_{im}$（$m = 1, \ldots, N$），可以把表示 (16.18) 推广到 $\mathrm{SU}(N)$：

\begin{equation*}
S ^ {mm '} = c _ {m} ^ {\dagger} c _ {m '} ,
\tag{16.21}
\end{equation*}

它们满足 $\mathrm{SU}(N)$ 代数 (16.9)。费米子的约束推广为

\begin{equation*}
\sum_ {m = 1} ^ {N} c _ {im} ^ {\dagger} c _ {im} = NS ,
\tag{16.22}
\end{equation*}

$S$ 是广义“自旋大小”。键算符定义为

\begin{equation*}
\mathcal {D} _ {ij} = \sum_ {m} c _ {im} ^ {\dagger} c _ {jm}.
\tag{16.23}
\end{equation*}

$SU(N)$ 反铁磁费米子（AFM-F）Heisenberg 模型由 (16.20) 推广：

\begin{align*}
\mathcal {H} ^ {\mathrm{AFM-F}} ( N ) & = \frac {J}{N} \sum_{\langle ij \rangle} \sum_ {mm '} S _ {i} ^ {mm '} S _ {j} ^ {m'm}\\
& = - \frac {J}{N} \sum_{\langle ij \rangle} : \mathcal {D} _ {ij} ^ {\dagger} \mathcal {D} _ {ij} : .
\tag{16.24}
\end{align*}

\section{生成泛函}

生成哈密顿量定义为

\begin{equation*}
\mathcal {H} [ j ] = \mathcal {H} - \sum_ {imm '} j _ {imm '} ( \tau ) a _ {im} ^ {\dagger} a _ {im '} ,
\tag{16.25}
\end{equation*}

$\tau\in[0,\beta)$ 是虚时间。约束 (16.6) 或 (16.22) 由与 $\mathcal{H}[j]$ 对易的投影算符 $P_{S}$ 强制。虚时生成泛函（见 (D.1)）为

\begin{align*}
Z [ j ] & = \mathrm{Tr} \, P _ {S} T _ {\tau} \left[ \exp \left( - \int_ {0} ^ {\beta} d \tau \, \mathcal {H} [ j ] \right) \right]\\
& = \lim_{\epsilon \to 0} \mathrm{Tr} \, T _ {\tau} \prod_{\tau_ {n} = \epsilon} ^ {\beta} \left[ P _ {S} ( \tau ) \exp \left( - \epsilon \mathcal {H} [ j ( \tau_ {n} ) ] \right) \right] .
\tag{16.26}
\end{align*}

用约束的积分表示

\begin{equation*}
P _ {S} ( \tau ) = \int \mathcal {D} \lambda \, \exp \left[ - i \epsilon \sum_ {im} \lambda_ {i} ( \tau ) \left( a _ {im} ^ {\dagger} a _ {im} - S \right) \right] ,
\tag{16.27}
\end{equation*}

约束场的测度为

\begin{equation*}
\int \mathcal {D} \lambda = \lim_{\epsilon \rightarrow 0} \prod_{i \tau} \epsilon \int_{- \pi / \epsilon} ^ {\pi / \epsilon} d \lambda_{i \tau}.
\tag{16.28}
\end{equation*}

(16.27) 的指数与 $\mathcal{H}[j]$ 可以合并进一个指数，因为它们对易。

循附录 D，为生成泛函构造相干态路径积分，对 Schwinger 玻色子与受约束费米子哈密顿量使用统一记号：

\begin{align*}
Z [ j ] & = \int_{- \infty} ^ {\infty} \mathcal {D} \lambda \int \mathcal {D} ^ {2} \mathbf {z} \; \exp \Bigg\{ - \int_ {0} ^ {\beta} d \tau \Big[ \sum_ {im} z _ {im} ^ {*} \partial_ {\tau} z _ {im} + H [ j ]\\
& \quad + i \sum_ {im} \lambda_ {i} ( \tau ) \left( z _ {im} ^ {*} z _ {im} - S \right) \Big] \Bigg\} ,
\tag{16.29}
\end{align*}

$z$ 对玻色子是复变量，对费米子是 Grassmann 变量。哈密顿量函数为

\begin{equation*}
H [ j ] = - \frac {J}{N} \sum_{\langle ij \rangle} \mathcal {Z} _ {ij} ^ {*} \mathcal {Z} _ {ij} - \sum_ {imm '} j _ {imm '} ( \tau ) z _ {im} ^ {*} z _ {im '} ,
\tag{16.30}
\end{equation*}

其中

\begin{equation*}
\mathcal {Z} = \left\{ \begin{array}{ll} \sum_ {m} z _ {im} ^ {*} z _ {jm} & \mathrm{FM\mbox{-}B} \\ \sum_ {m} z _ {im} z _ {jm} & \mathrm{AFM\mbox{-}B} \\ \sum_ {m} z _ {im} ^ {*} z _ {jm} & \mathrm{AFM\mbox{-}F} \end{array} \right. ,
\tag{16.31}
\end{equation*}

FM-B、AFM-B 与 AFM-F 哈密顿量分别已在 (16.7)、(16.15) 与 (16.24) 定义。

\section{Hubbard--Stratonovich 变换}

Lagrange 量中的双二次（四变量）项无法在路径积分中直接积分。用 Hubbard--Stratonovich 恒等式

\begin{align*}
& \exp \left[ \frac {J}{N} \mathcal {Z} ^ {*} \mathcal {Z} \epsilon \right] = \int_{- \infty} ^ {\infty} d ^ {2} Q\\
& \quad \times \exp \left[ - \left( \mathcal {Z} ^ {*} Q + \mathcal {Z} Q ^ {*} + N \frac {| Q | ^ {2}}{J} \right) \epsilon \right] ,
\tag{16.32}
\end{align*}

其中

\begin{equation*}
d ^ {2} Q \equiv \frac {\epsilon N}{\pi J} \, \mathrm{Re} Q \, d \mathrm{Im} Q .
\tag{16.33}
\end{equation*}

对每个时间步的每条键 $\langle ij\rangle$ 引入复变量 $Q_{ij}(\tau)$，积分测度定义为

\begin{equation*}
\mathcal {D} ^ {2} Q \equiv \prod_{\langle ij \rangle \tau} d ^ {2} Q _ {ij} ( \tau ).
\tag{16.34}
\end{equation*}

用 (16.27)、(16.29) 与 (16.32)：

\begin{align*}
Z [ j ] & = \int \mathcal {D} ^ {2} Q \, \mathcal {D} \lambda \, \mathcal {D} ^ {2} z\\
& \quad \times \exp \left[ - \int_ {0} ^ {\beta} d \tau \left( L [ j ] + N \frac {| Q _ {ij} | ^ {2}}{J} - iNS \sum_ {i} \lambda_ {i} \right) \right] ,
\tag{16.35}
\end{align*}

Lagrange 量对 $z$ 变量是二次的：

\begin{align*}
L [ \mathbf {z} ^ {*}, \mathbf {z}, j ] & = \sum_ {im} z _ {im} ^ {*} \partial_ {\tau} z _ {im} + \sum_{\langle ij \rangle} \left( \mathcal {Z} _ {ij} ^ {\dagger} Q _ {ij} + \mathcal {Z} _ {ij} Q _ {ij} ^ {*} \right)\\
& \quad - \sum_ {imm '} \left( j _ {imm '} + i \lambda_ {i} \delta_{mm '} \right) z _ {im} ^ {*} z _ {im '} .
\tag{16.36}
\end{align*}

对只含 $z^{*}z$ 项的正规 Lagrange 量（如 FM-B 与 AFM-F 情形），Green 函数矩阵 $\hat{G}$ 定义为

\begin{align*}
L & = \mathbf {z} ^ {*} \hat {G} ^ {- 1} [ j ] \mathbf {z} ,\\
\hat {G} ^ {- 1} & = \partial_ {\tau} + \hat {\lambda} + \hat {Q} - \hat {j} ,
\tag{16.37}
\end{align*}

$\partial_{\tau}$ 已在 (D.9) 定义，$\hat{\lambda},\hat{Q},\hat{j}$ 由 (16.36) 读出。

对反常 Lagrange 量（如 AFM-B）：

\begin{align*}
\mathcal {L} & = \left( \mathbf {z} _ {A} ^ {*}, \mathbf {z} _ {B} \right) \hat {G} ^ {- 1} [ j ] \binom{\mathbf {z} _ {A}}{\mathbf {z} _ {B} ^ {*}} ,\\
\hat {G} ^ {- 1} & = \left( \begin{array}{cc} \partial_ {\tau} + \hat {\lambda} - \hat {j} & \hat {Q} \\ \hat {Q} ^ {\dagger} & - \partial_ {\tau} + \hat {\lambda} - \hat {j} \end{array} \right) ,
\tag{16.38}
\end{align*}

$z_{A}, z_{B}$ 分别是子晶格 A、B 的变量。

循 (C.13) 与 (D.11) 积掉 $z$ 变量，得到辅助场的路径积分：

\begin{align*}
Z [ j ] & = \int \mathcal {D} ^ {2} Q \, \mathcal {D} \lambda \, \exp \left[ - N \mathcal {S} [ \lambda, Q, j ] \right]\\
\mathcal {S} & = - \frac {\eta}{N} \mathrm{Tr}_{\tau im} \left( \ln \hat {G} [ j ] \right) + \int_ {0} ^ {\beta} d \tau \left( \sum_{\langle ij \rangle} \frac {| Q _ {ij} | ^ {2}}{J} - iS \sum_ {i} \lambda_ {i} \right) .
\tag{16.39}
\end{align*}

这个表达式是以 $N$ 为大参数的最陡下降展开的出发点。大 $N$ 展开将在第 17、18 章综述。

\section{关联函数}

连通的两自旋关联函数为

\begin{align*}
\tilde {R} ^ {mm '} ( 1, 2 ) & = \left\langle a _ {m, 1} ^ {\dagger} a _ {m ', 1} a _ {m ', 2} ^ {\dagger} a _ {m, 2} \right\rangle\\
& = Z ^ {- 1} \left. \frac {\delta^ {2} \ln Z}{\delta j_{mm', 1} \, \delta j_{m'm, 2}} \right| _ {j = 0}\\
& = Z ^ {- 1} \int \mathcal {D} ^ {2} Q \, \mathcal {D} \lambda \Big[ - N \frac {\delta^ {2} \mathcal {S}}{\delta j_{mm', 1} \, \delta j_{m'm, 2}}\\
& \quad + N ^ {2} \frac {\delta \mathcal {S}}{\delta j_{mm', 1}} \frac {\delta \mathcal {S}}{\delta j_{m'm, 2}} \exp \left( - N \mathcal {S} [ \lambda, Q ] \right) \Big] ,
\tag{16.40}
\end{align*}

1、2 是时空中的点，被积函数中所有泛函都在 $j=0$ 处取值。由 $\mathcal{S}$ 对味指标的对称性，第二项正比于 $\delta_{mm'}$，$\tilde{R}^{m\neq m'}$ 与 $m,m'$ 无关。对 $SU(2)$，FM-B 与 AFM-F 关联与通常自旋关联的关系为

\begin{equation*}
\tilde {R} _ {N = 2} ^ {m \neq m '} ( 1, 2 ) = \left\langle S _ {i _ {1}} ^ {+} ( \tau_ {1} ) S _ {i _ {2}} ^ {-} ( \tau_ {2} ) \right\rangle .
\tag{16.41}
\end{equation*}

AFM-B 情形，子晶格转动 (16.11) 蕴含

\begin{equation*}
\tilde {R} _ {N = 2} ^ {m \neq m '} ( 1, 2 ) = \left\{ \begin{array}{ll} \left\langle S _ {i _ {1}} ^ {+} ( \tau_ {1} ) S _ {i _ {2}} ^ {-} ( \tau_ {2} ) \right\rangle & i _ {1}, i _ {2} \in A \\ - \left\langle S _ {i _ {1}} ^ {+} ( \tau_ {1} ) S _ {i _ {2}} ^ {+} ( \tau_ {2} ) \right\rangle & i _ {1} \in A, i _ {2} \in B \end{array} \right. .
\tag{16.42}
\end{equation*}

\begin{chapbib}
\item 大 $N$ 方法起源于场论，见第 14 章及其文献注释。
\item 相互作用电子的首次大 $N$ 展开应用于 Coqblin--Schrieffer 模型（Kondo 模型的大 $N$ 推广）：
  N. Read and D. Newns, J. Phys. C \textbf{16}, 3273 (1983)。
\item Coleman 把大 $N$ 展开用于 Anderson 模型的从属玻色子表示：
  P. Coleman, Phys. Rev. B \textbf{35}, 5072 (1987)。
\item 用 $1/N$ 展开导出重费米子的 Fermi 液体参数：
  A. Auerbach and K. Levin, Phys. Rev. Lett. \textbf{57}, 877 (1986)；
  A. Millis and P. Lee, Phys. Rev. B \textbf{35}, 3394 (1987)。
\item 大 $N$ 近似在重费米子与铜氧化物超导体中应用的详尽综述：
  A. Hewson, \textit{The Kondo Problem to Heavy Fermions} (Cambridge, 1993)；
  N. E. Bickers, Rev. Mod. Phys. \textbf{59}, 845 (1987)；
  G. Kotliar, \textit{Correlated Electron Systems}, edited by V. J. Emery (World Scientific, 1993)；
  K. Levin, J. H. Kim, J. P. Lu, and Q. Si, Physica C \textbf{175}, 449 (1991)。
\end{chapbib}
